\section{Two-body problem}

Having dealt with the original generalisation of the Kepler problem
proposed in \cite{1}, a natural question arises about the two-body
problem. In the classical Kepler problem, there is no fundamental
difference between one and two bodies: the latter still leads to the
Kepler problem for a single body of reduced mass, revolving around the
center of mass. The reduction is possible due to the symmetries of the Euclidean space, which generate boosts, and correspond closely to the motion: relative positions of two particles follow a geodesic. This is not the case for the Heisenberg group, where the group operation \eqref{groupop} does not preserve geodesics, as discussed in~detail by the authors of~\cite{1}~-- the difference leads them to pose the integrability question also for the two-body case. In what follows, we give a decisive answer: the two-body problem on the Heisenberg group is not integrable.

For two point masses $m_1$ and $m_2$, whose positions are group
elements $g_k = (x_k,y_k,z_k)$ we will take the Hamiltonian to be
\begin{gather}
 H = \frac{1}{2m_1}\bigg(\bigg(p_{x_1}-\frac12 y_1 p_{z_1}\bigg)^2
 + \bigg(p_{y_1}+\frac12 x_1 p_{z_1}\bigg)^2\bigg)\nonumber
 \\ \hphantom{H =}
 {}+ \frac{1}{2m_2}\bigg(\bigg(p_{x_2}-\frac12 y_2 p_{z_2}\bigg)^2 +
 \bigg(p_{y_2}+\frac12 x_2 p_{z_2}\bigg)^2\bigg) +
 V\big(\rho\big(g_1^{-1}\cdot g_2\big)\big),\label{eq:2b}
\end{gather}
where the potential is specified by
\begin{equation*}
 V(\rho) = -\frac{\kappa m_1 m_2}{\rho},\qquad
 \rho(g) = \sqrt{\big(x^2+y^2\big)^2+{16}z^2}.
\end{equation*}

We note that the following first integrals are ``known'':
\begin{alignat*}{3}
& I_1 = p_{x_1} +\frac12 y_1 p_{z_1} + p_{x_2} +\frac12 y_2 p_{z_2},
 \qquad&&
 I_2 = p_{y_1} -\frac12 x_1 p_{z_1} + p_{y_2} -\frac12 x_2 p_{z_2},&
 \\
 &I_3 = p_{z_1}+p_{z_2}=\{I_1,I_2\},
 \qquad&&
 I_4 = y_1 p_{x_1}-x_1 p_{y_1} + y_2 p_{x_2} - x_2 p_{y_2},&
\end{alignat*}
and they satisfy
\begin{equation*}
 \{I_1,I_4\}=I_2, \qquad
 \{I_2,I_4\} = -I_1.
\end{equation*}
Additionally, as for the Kepler--Heisenberg problem,
\begin{equation*}
 J = x_1 p_{x_1} + y_1 p_{y_1} + 2 z_1 p_{z_1} + x_2 p_{x_2} + y_2 p_{y_2} + 2 z_2 p_{z_2},
\end{equation*}
is such that $\dot{J}=2H$. That is, we have 5 first integrals,
although they do not all commute, and on the zero-energy level $J$
becomes the sixth integral. The question, as before, is whether there
exist enough (here: six) commuting integrals.

Now, we make linear canonical transformation
\begin{alignat*}{3}
 &u_1 = \frac{1}{\sqrt{2}}(y_1-\mathrm{i} x_1),\qquad&&
 p_{u_1} = \frac{1}{\sqrt{2}}(p_{y_1}+\mathrm{i}p_{x_1}),&
 \\
& v_1 = \frac{1}{\sqrt{2}}(y_1+\mathrm{i} x_1),\qquad&&
 p_{v_1} =\frac{1}{\sqrt{2}}(p_{y_1}-\mathrm{i}p_{x_1}),&
 \\
& w_1 = z_1 +z_2, \qquad &&
 p_{w_1} = \frac12(p_{z_1}+p_{z_2}),&
 \\
 &u_2 = \frac{1}{\sqrt{2}}(y_2-\mathrm{i} x_2),\qquad&&
 p_{u_2} = \frac{1}{\sqrt{2}}(p_{y_2}+\mathrm{i}p_{x_2}),&
 \\
 &v_2 = \frac{1}{\sqrt{2}}(y_2+\mathrm{i} x_2),\qquad&&
 p_{v_2} =\frac{1}{\sqrt{2}}(p_{y_2}-\mathrm{i}p_{x_2}),&
 \\
& w_2 = z_1 -z_2, \qquad&&
 p_{w_2} = \frac12(p_{z_1}-p_{z_2}).&
\end{alignat*}
In the new variables the Hamiltonian reads
\begin{gather*}
 H = \frac{((p_{w_1}+p_{w_2})u_1-2\mathrm{i}p_{v_1})
 ((p_{w_1}+p_{w_2})v_1+2\mathrm{i}p_{u_1})}{4m_1}
 \\ \hphantom{ H =}
{} + \frac{((p_{w_1}-p_{w_2})u_2-2\mathrm{i}p_{v_2})
 ((p_{w_1}-p_{w_2})v_2+2\mathrm{i}p_{u_2})}{4m_2} -\frac{\kappa m_1 m_2}{\rho_{12}},
\end{gather*}
where
\begin{equation*}
 \rho_{12} = 2\sqrt{(u_1-u_2)^2(v_1-v_2)^2-(v_1 u_2 - v_2 u_1 + 2\mathrm{i}w_2)^2}.
\end{equation*}

The particular solution is almost as before
\begin{equation*}
 p_{w_1}=p_{w_1}(0),\qquad
 p_{w_2} = -2at,\qquad
 a = \frac{m_1 m_2\kappa}{8w_2|w_2|},\qquad
 w_2=w_2(0), \qquad
 w_1=w_1(0),
\end{equation*}
and all other phase variables equal to zero. The solution must not be constant, so $w_2(0)\neq 0$, but other parameters are not restricted.

Linear variations of the variables
$(u_1, p_{v_1},u_2,p_{v_2},v_1,p_{u_1},v_2,p_{u_2},w_1, w_2,
p_{w_1},p_{w_2})$, which we will denote by $\xi=(\xi_1, \ldots, \xi_{12})$, then
satisfy the variational equations
\begin{equation}
 \label{eq:varg}
 \dot \xi = A\xi, \qquad A = \begin{bmatrix}
 A_1 & 0 & 0\\
 0 & A_2 & 0\\
 0 & 0 & A_3
 \end{bmatrix}\!,
\end{equation}
where
\begin{equation*}
 A_1 = \begin{bmatrix}
 \tau -\tau_0& 1 & 0 & 0\\
 (\tau-\tau_0)^2 & \tau-\tau_0 & -1 & 0\\
 0 & 0 & -\mu(\tau+\tau_0) & \mu \\
 1 & 0 & \mu(\tau+\tau_0)^2 & -\mu(\tau+\tau_0)
 \end{bmatrix}\!,\qquad
 A_3 = \begin{bmatrix}
 0 & 0 & 0 & 0\\
 0 & 0 & 0 & 0\\
 0 & 0 & 0 & 0 \\
 0 & 4\mathrm{i}/w_2 & 0 & 0
 \end{bmatrix}\!,
\end{equation*}
and
\begin{equation*}
 A_2 = \begin{bmatrix}
 \tau_0-\tau & 1 & 0 & 0\\
 (\tau_0-\tau)^2 & \tau_0-\tau & 1 & 0\\
 0 & 0 & \mu(\tau+\tau_0) & \mu \\
 -1 & 0 & \mu(\tau+\tau_0)^2 & \mu(\tau+\tau_0)
 \end{bmatrix}\!.
\end{equation*}
To obtain the above form we use the following rescalings
\begin{equation*}
 t = m_1\tau,\qquad
 a = \frac{\mathrm{i}}{m_1},\qquad
 \mu = \frac{m_1}{m_2}, \qquad
 p_{w_1}=2 \mathrm{i} \tau_0.
\end{equation*}

\begin{Theorem} %\label{thm:1}
 If $\mu\neq -1$ then the two-body problem on the Heisenberg group is
 not integrable in the Liouville sense.
\end{Theorem}
\begin{proof}
 If the system generated by~\eqref{eq:2b} is integrable then by
 Theorem~\ref{thm:MoRa}, the identity component of differential
 Galois group of variational equations~\eqref{eq:varg} is Abelian.
 This implies that the same property is shared by the differential Galois groups
 of the subsystems of~\eqref{eq:varg}, which have the form
 $\dot\eta=A_i \eta$, $\eta\in\mathbb{C}^4$, for $i=1,2,3$. We
 consider the first of them. It has particular solution
 \begin{equation}
 \label{par}
 \eta(\tau)=(1, \tau_0-\tau,1, \tau_0+\tau).
 \end{equation}
 Using the d'Alambert method, see~\cite{Walter}, we can reduce the
 dimension of the system by one. But~assuming that $\tau_0=0$ we
 achieve more. Namely, linear transformation $\eta\mapsto Q\eta$ with
 $Q=Q(\tau)$ given by
 \begin{equation*}
 Q = \begin{bmatrix}
 1 & 0 & 0 & 0\\
 -\tau & 1 & 0 & 0\\
 1 & 2\tau & -1 & 0\\
 \tau & -1-2\tau^2 & \tau & -\frac{1}{\mu}
 \end{bmatrix}\!,
 \end{equation*}
 brings it to the form $\dot{\eta}=\widetilde{A}_1\eta$, where
 \begin{equation*}
 \widetilde{A}_1 =Q^{-1}\bigg(A_1 Q -\frac{\mathrm{d}}{\mathrm{d}\tau}Q\bigg) = \begin{bmatrix}
 0 & 1 & 0 & 0\\
 0 & 0 & 1 & 0\\
 0 & p_2 & p_1 & 1\\
 0 & 0 & 0 & 0
 \end{bmatrix}\!,\qquad p_1 = 2(1-\mu)\tau,\qquad p_2 =
 3+\mu+4\mu\tau^2.
 \end{equation*}
 Thus the transformed system has block-triangular structure which is
 quite simple: the first coordinate does not enter, while the fourth
 is constant. In other words, to obtain a particular solution, it is
 enough to assume $\eta_4=0$, and choose the subsystem corresponding to the second and third
 components:
 \begin{gather*}
 \eta_2'(\tau) = \eta_3(\tau),
 \\
 \eta_3'(\tau) = p_2(\tau) \eta_2(\tau) + p_1(\tau)\eta_3(\tau).
 \end{gather*}
 As a single equation it reads
 \begin{equation*}
 \eta_2'' = 2(1-\mu)\tau \eta_2' + \big(3+\mu+4\mu \tau^2\big)\eta_2,
 \end{equation*}
 which, after the change $\eta_2=\exp\big[(1-\mu)\tau^2/2\big]w(\tau)$, becomes
 \begin{equation}
 \label{eq:w2b}
 w''(\tau) -(1+\mu)\big[2 +(1+\mu)\tau^2\big]w(\tau)=0.
 \end{equation}
 It is, again, the parabolic cylinder equation~\eqref{eq:5} with parameters
 \begin{equation*}
 \alpha^2= (1+\mu)^2, \qquad
 \beta=0, \qquad
 \gamma = 2(1+\mu).
 \end{equation*}
 Let us assume that $\mu\neq-1$. Then $\alpha\neq 0$ and
 \[
 \frac{\beta^2-\gamma}{\alpha}=-2\sgn(1+\mu)
 \]
 is not an odd integer. Hence, by Theorem~\ref{thm:par} the
 differential Galois group of equation~\eqref{eq:w2b} is~$\mathrm{SL}(2,\mathbb{C})$. This ends the proof.
\end{proof}
